\usepackage{geometry}
\geometry{verbose,tmargin=1in,bmargin=1in,lmargin=1in,rmargin=1in}

\usepackage{amsmath, amssymb, amsfonts, bm, mathtools}
\usepackage{amsthm}
\usepackage{xcolor}
\definecolor{darkblue}{rgb}{0,0,.5}
\usepackage{graphicx}
\usepackage{subfigure}

\newcommand{\theHalgorithm}{\arabic{algorithm}}

\usepackage[numbers]{natbib}

\usepackage[colorlinks=true,allcolors=darkblue]{hyperref}       % hyperlinks
\usepackage{url}            % simple URL typesetting
\usepackage{booktabs}       % professional-quality tables
\usepackage{amsfonts}       % blackboard math symbols
\usepackage{nicefrac}       % compact symbols for 1/2, etc.
\usepackage{microtype}      % microtypography

\allowdisplaybreaks

\usepackage{float}
\usepackage{multirow}
\usepackage{footnote}
\usepackage{dsfont}
% \usepackage{mathabx}

\usepackage{algorithm}
\usepackage{algorithmic}
\usepackage{nicefrac}


\usepackage{hyperref}
\usepackage[capitalize]{cleveref}
\usepackage{crossreftools}


\usepackage{tikz}
\usepackage{overpic}


\definecolor{cc}{RGB}{125,0,0}
\definecolor{jd}{RGB}{0,125,0}
\definecolor{ha}{RGB}{0,0,200}

\newcommand{\cc}[1]{\textcolor{cc}{[CC: #1]}}
\newcommand{\jd}[1]{\textcolor{jd}{[JD: #1]}}
\newcommand{\ha}[1]{\textcolor{ha}{[HA: #1]}}

 % \usepackage[inline]{showlabels}


% \newcommand{\Pr}{\mathsf{Pr}}

\usepackage{dsfont}


\newcommand{\ind}{\mathds{1}}
\newcommand{\var}{\mathsf{Var}}
\newcommand{\cov}{\mathsf{Cov}}
\newcommand{\Ep}{\mathbb{E}}
\newcommand{\E}{\mathbb{E}}
\newcommand{\Prb}{\mathbb{P}}
\newcommand{\R}{\mathbb{R}} % real numbers
\newcommand{\N}{\mathbb{N}} % natural numbers
\newcommand{\Z}{\mathbb{Z}} % integers
\renewcommand{\P}{\mathbb{P}}
\newcommand{\mc}[1]{\mathcal{#1}}
\newcommand{\brc}[1]{\left\{{#1}\right\}}
\newcommand{\prn}[1]{\left({#1}\right)} % parentheses
\newcommand{\prnbig}[1]{\big({#1}\big)} % parentheses
\newcommand{\prns}[1]{({#1})} % parentheses small
\newcommand{\prnBig}[1]{\Big({#1}\Big)} % parentheses
\newcommand{\prnBigg}[1]{\Bigg({#1}\Bigg)} % parentheses
\newcommand{\brk}[1]{\left[{#1}\right]} % bracket
\newcommand{\brkbig}[1]{\big[{#1}\big]} % parentheses
\newcommand{\brkBig}[1]{\Big[{#1}\Big]} % parentheses
\newcommand{\brkBigg}[1]{\Bigg[{#1}\Bigg]} % parentheses
\newcommand{\norm}[1]{\left\|{#1}\right\|} % norm
\newcommand{\normbig}[1]{\big\|{#1}\big\|} % norm
\newcommand{\abs}[1]{\left|{#1}\right|} % norm
\newcommand{\what}[1]{\widehat{#1}}
\newcommand{\wt}[1]{\widetilde{#1}}
\newcommand{\wb}[1]{\overline{#1}}
\newcommand{\wbmj}[1]{{#1}^+}
\newcommand{\majY}{Y^+}
\newcommand{\majmark}[1]{{#1}^+}
\newcommand{\sgn}{\mathsf{sign}}
\newcommand{\pred}{\mathsf{d}}
\newcommand{\dist}{\mathsf{dist}}
\newcommand{\NN}{\mathsf{NN}}
\newcommand{\normal}{\mathsf{N}}
\newcommand{\bindist}{\mathsf{Binomial}}
\newcommand{\maj}{\mathsf{maj}}
\newcommand{\dstr}{\mathcal{P}}
\newcommand{\lr}{^{\textup{lr}}}
\newcommand{\sem}{^{\textup{sp}}}
\newcommand{\cs}{^{\textup{cs}}}
\newcommand{\mv}{^{\textup{mv}}}
\newcommand{\ds}{^{\textsc{ds}}}
\newcommand{\dsp}{^{\textsc{ds-p}}}
\newcommand{\glad}{^{\textsc{glad}}}
\newcommand{\gladp}{^{\textsc{glad-p}}}
\newcommand{\amv}{^{\textup{amv}}}
\newcommand{\mm}{^{\textup{mm}}}
\newcommand{\proj}{\mathsf{P}}
\newcommand{\proju}{\proj_{u\opt}}
\newcommand{\projperp}{\proj_{u\opt}^\perp}
\newcommand{\Flink}{\mathcal{F}_{\mathsf{link}}}
\newcommand{\lipconst}{\mathsf{L}}
\newcommand{\loss}{\ell}
\newcommand{\poploss}{L}
\newcommand{\popsignloss}{L^{\textup{sgn}}}
\newcommand{\flink}{\mathcal{F}_{\mathsf{link}}}
\newcommand{\fsymlink}{\mathcal{F}_{\mathsf{link}}^0}
\newcommand{\fsymlinkL}{\mathcal{F}_{\mathsf{link}}^{\mathsf{L}}}
\newcommand{\fsymlinksp}{\mathcal{F}_{\mathsf{link}}^{\mathsf{sp}}}
% Norms for convenience, along with sizes
\newcommand{\lone}[1]{\norm{#1}_1} % l1 norm
\newcommand{\ltwo}[1]{\norm{#1}_2} % l2 norm
\newcommand{\linf}[1]{\norm{#1}_\infty} % l-infinity norm
\newcommand{\lzero}[1]{\norm{#1}_0} % l0 norm
\newcommand{\dnorm}[1]{\norm{#1}_*} % Dual norm
\newcommand{\lfro}[1]{\left\|{#1}\right\|_{\rm Fr}} % Frobenius norm
\newcommand{\ltv}[2]{\mathsf{TV}\prn{#1, #2}}
\newcommand{\matrixnorm}[1]{\left|\!\left|\!\left|{#1}
	\right|\!\right|\!\right|} % Matrix norm with three bars
\newcommand{\normbigg}[1]{\bigg\|{#1}\bigg\|} % A norm with 1 argument and bigg
% brackets.
\newcommand{\dnormbigg}[1]{\bigg\|{#1}\bigg\|_*} % A norm with 1 argument and bigg
% brackets.
\newcommand{\lonebigg}[1]{\normbigg{#1}_1} % l1 norm
\newcommand{\ltwobigg}[1]{\normbigg{#1}_2} % l2 norm
\newcommand{\ltwobig}[1]{\normbig{#1}_2} % l2 norm
\newcommand{\ltwoBig}[1]{\normBig{#1}_2} % l2 norm
\newcommand{\linfbigg}[1]{\normbigg{#1}_\infty} % l-infinity norm
\newcommand{\norms}[1]{\|{#1}\|} % A norm with 1 argument and normal (small)
% brackets.
\newcommand{\matrixnorms}[1]{|\!|\!|{#1}|\!|\!|} % Small matrix norm
\newcommand{\matrixnormbigg}[1]{\bigg|\!\bigg|\!\bigg|{#1}\bigg|\!\bigg|\!\bigg|} % Small matrix norm
\newcommand{\lones}[1]{\norms{#1}_1} % l1 norm with small brackets
\newcommand{\ltwos}[1]{\norms{#1}_2} % l2 norm with small brackets
\newcommand{\linfs}[1]{\norms{#1}_\infty} % l-infinity norm with small brackets
\newcommand{\lzeros}[1]{\norms{#1}_0} % l0 norm with small brackets

\newcommand{\opnorm}[1]{\matrixnorm{#1}_{\rm op}} % Operator norm with three bars
\newcommand{\opnormbigg}[1]{\matrixnormbigg{#1}_{\rm op}} % Operator norm with three bars
\newcommand{\opnorms}[1]{\matrixnorms{#1}_{\rm op}} % Small operator norm with three bars
\newcommand{\<}{\langle} % Angle brackets
\renewcommand{\>}{\rangle}

\DeclareMathOperator*{\argmax}{argmax}
\DeclareMathOperator*{\argmin}{argmin}

\newcommand{\sX}{\mathcal{X}} % input space
\newcommand{\sY}{\mathcal{Y}} % output space


\newcommand{\simiid}{\stackrel{\textup{iid}}{\sim}}
\newcommand{\indic}[1]{1\!\left\{{#1}\right\}}

\newcommand{\cd}{\rightsquigarrow}
% \newcommand{\cd}{\stackrel{d}{\rightarrow}}
\newcommand{\cp}{\stackrel{p}{\rightarrow}}
\newcommand{\cas}{\stackrel{a.s.}{\rightarrow}}

\makeatletter
\long\def\@makecaption#1#2{
  \vskip 0.8ex
  \setbox\@tempboxa\hbox{\small {\bf #1:} #2}
  \parindent 1.5em  %% How can we use the global value of this???
  \dimen0=\hsize
  \advance\dimen0 by -3em
  \ifdim \wd\@tempboxa >\dimen0
  \hbox to \hsize{
    \parindent 0em
    \hfil 
    \parbox{\dimen0}{\def\baselinestretch{0.96}\small
      {\bf #1.} #2
      %%\unhbox\@tempboxa
    } 
    \hfil}
  \else \hbox to \hsize{\hfil \box\@tempboxa \hfil}
  \fi
}
\makeatother

\newcommand{\defeq}{\coloneqq}
\newcommand{\Ltwo}[1]{\norm{#1}_{L^2}}
\newcommand{\LtwoP}[1]{\norm{#1}_{L^2(\Prb)}}
\newcommand{\LtwoPbold}[1]{\norm{#1}_{L^2(\Prb)}}
\newcommand{\LtwoPn}[1]{\norm{#1}_{L^2(\Prb_n)}}
\newcommand{\LtwoPns}[1]{\norms{#1}_{L^2(\Prb_n)}}
\newcommand{\LtwoQ}[1]{\norm{#1}_{L^2(Q)}}
\newcommand{\LtwoQn}[1]{\norm{#1}_{L^2(Q_n)}}
\newcommand{\LtwoQns}[1]{\norms{#1}_{L^2(Q_n)}}
\newcommand{\lipnorm}[1]{\norm{#1}_{\textup{Lip}}}

\newcommand{\half}{\frac{1}{2}}
\newcommand{\event}{\mc{E}}

\newcommand{\conv}{\textup{Conv}}
\newcommand{\card}{\textup{card}}
\newcommand{\starhull}{\textup{Star}}
\newcommand{\hinge}[1]{\left({#1}\right)_+}
\newcommand{\hinges}[1]{({#1})_+}

\newcommand{\red}[1]{\textcolor{red}{#1}}
\newcommand{\blue}[1]{\textcolor{blue}{#1}}
\newcommand{\darkblue}[1]{\textcolor{darkblue}{#1}}

\definecolor{innerboxcolor}{rgb}{.9,.95,1}
\definecolor{outerlinecolor}{rgb}{.6,0,.2}

\newcommand{\jcdcomment}[1]{\fcolorbox{outerlinecolor}{innerboxcolor}{
    \begin{minipage}{.9\textwidth}
      \red{\bf JCD Comment:} {#1}
  \end{minipage}} \\
}

\newcommand{\opt}{^\star}
\newcommand{\subopt}{_\star}
\newcommand{\gme}{\what{\theta}_{n,m}} %General multilabel estimator
\newcommand{\tmle}{\what{\theta}\lr_{n,m}} %MLE estimator without aggregation (logistic regression)
\newcommand{\mve}{\what{\theta}\mv_{n,m}} %Majority vote estimator
\newcommand{\spe}{\what{\theta}\sem_{n,m}} %Semiparametric estimator
\newcommand{\csmle}{\what{\theta}\cs_{n,m}} %Crowdsourcing estimator
\newcommand{\dse}{\what{\theta}\ds_{n,m}} %Dawid-Skene estimator
\newcommand{\dspe}{\what{\theta}\dsp_{n,m}} %Dawid-Skene soft estimator
\newcommand{\glade}{\what{\theta}\glad_{n,m}} %GLAD estimator
\newcommand{\gladpe}{\what{\theta}\gladp_{n,m}} %GLAD soft estimator
\newcommand{\G}{\mathbb{G}}
